\documentclass[11pt, a4paper, oneside]{article}
\usepackage[ngerman]{babel}
\usepackage{worksheet}

\begin{document}
	\author{L. Bung}
	\title{Trigonometrische Funktionen}
	\subject{Mathematik}
	\class{FHR1}
	\maketitle
	
	\begin{minipage}{.75\textwidth}
		\task{Der Einheitskreis}
			
		a) Füllen Sie die folgende Tabelle mithilfe des Geogebra-Applets aus.\\
		
		\begin{tabular}{|p{1.5cm}|p{2cm}|p{1.5cm}|p{1.5cm}|p{1cm}|p{1cm}|}
			\hline
			Winkel $\alpha$ & Bogenlänge & $x$ & $y$ & $\sin(\alpha)$ & $\cos(\alpha)$\\
			\hline
			\hline
			$0^\circ$ &&&&&\\
			\hline
			$30^\circ$ &&&&&\\
			\hline
			& $\pi$ &&&&\\
			\hline
			&&& $-1$&&\\
			\hline
			& $2\pi$ &&&&\\
			\hline
		\end{tabular}
	\end{minipage}
	\begin{minipage}{.25\textwidth}
		\centering
		\qrlink{https://www.geogebra.org/m/ykmewjec}
	\end{minipage}
	
	b) Beschreiben Sie, was Ihnen auffällt.
	
	\lines[2cm]
	
	c) Erklären Sie, warum das so ist.
	
	\lines[3cm]
	
	d) Finden Sie eine Formel, mit der man den Winkel $\alpha$ in die Bogenlänge umrechnen kann.
	
	\checkered[3cm]
	
	
	\task{Die Graphen der Sinus- und Cosinusfunktion}
	
	a) Zeichnen Sie den Graphen von $f(x) = \sin(x)$ und $g(x) = \cos(x)$ für $x \in \left[0; 2\pi\right]$.
	Stellen Sie dafür zuerst eine Wertetabelle auf.
	
	\emph{Wichtig: Der Taschenrechner muss dafür auf Bogenmaß (Radian) eingestellt werden!}
	
	\checkered[8cm]
	
	b) Stellen Sie eine Vermutung auf, wie die Graphen von $f(x)$ und $g(x)$ für $x < 0$ und $x > 2\pi$ weitergehen.
	Überprüfen Sie Ihre Vermutung, indem Sie mit dem Taschenrechner weitere Werte ausrechnen.
	
	\lines[3cm]
	
	c) Beschreiben Sie, in welchem Zusammenhang die Graphen von $f(x)$ und $g(x)$ stehen.
	
	\lines[3cm]
	
	
	\task{Die Tangensfunktion}
	
	Im Geogebra-Applet sehen Sie zusätzlich zum Sinus und Cosinus auch noch den Graphen der Tangens-Funktion.
	
	\begin{minipage}{.75\textwidth}
		a) Beschreiben Sie den Verlauf des Graphen der Tangens-Funktion.
		
		\lines[3cm]
	\end{minipage}
	\begin{minipage}{.25\textwidth}
		\qrlink{https://www.geogebra.org/m/kcfhyegj}
	\end{minipage}
	
	b) Die Tangens-Funktion hat an bestimmten Stellen Definitionslücken.
	Finden Sie heraus, welche Stellen das sind.
	
	\lines[2cm]
	
	c) Begründen Sie mithilfe der Definition des Tangens, weswegen die Definitionslücken genau an diesen Stellen auftreten.
	
	\lines[3cm]
\end{document}